\section{Introduction}

Adult obesity worldwide has more than doubled since 1990: in 2022, 2.5 billion adults were overweight, 890 million of them living with obesity — about one in eight people in the world \cite{ref_who2024}. A healthy body mass matters because excess body fat raises the risk of type 2 diabetes, cardiovascular disease, and several cancers \cite{ref_who2024}.

Yet a systematic view of human metabolism is missing \cite{ref_foresight2007}. Weight-loss advice usually acts on a single variable, calories in \cite{ref_hall2011}, while the body responds to many at once: activity \cite{ref_hall2012}, sleep \cite{ref_nedeltcheva2010}, food composition \cite{ref_hall2017}, and its own metabolic adaptation \cite{ref_fothergill2016}. The human metabolic system is multi-input and multi-output, so no single variable can control it \cite{ref_hall2022}; building on dynamic and control-engineering models of body mass \cite{ref_chow2008}, \cite{ref_navarro2011}, it needs a multi-input multi-output (MIMO) approach that accounts for all of its variables together, which is the model this paper develops.
